% GENERATED FILE. Edit canonical lessons, then run npm run generate:books.
% canonical-source-hash: fnv1a64:43c340097d7e5fef
% canonical-lessons: HI-C280-lagna, HI-C280-aisa, HI-W280-letter-ai, HI-C280-mujhelagtahai, HI-C280-kivajahse, HI-R280-opinion-and-reason

\chapter{What I Think, and Why}
\label{ch:hi-what-i-think-and-why}
\hlchaptermodality{280}

\begin{chapteropening}
\textbf{By the end of this chapter:} \emph{I can say what I think with mujhe lagtā hai ki and merī rāy meṁ, give a reason with kī vajah se and ke kāraṇ, and read and write \hi{ऐ} in aisā.}

Complete HI-R280-opinion-and-reason: say what you think and give a reason for it.
\end{chapteropening}

\section[\texorpdfstring{\hi{लगना} (lagnā)}{लगना (lagnā)}]{\hi{लगना} (lagnā) — to seem, to be felt}

\label{lesson:HI-C280-lagna}



\begin{quote}
Before the new one: say \textquotedblleft{}nice to meet you\textquotedblright{} the way you say it as you part. (\textbf{\hi{मिलकर} \hi{अच्छा} \hi{लगा}}, \emph{milkar acchā lagā}.) Then \textquotedblleft{}I have been hurt\textquotedblright{}. (\textbf{\hi{मुझे} \hi{चोट} \hi{लगी} \hi{है}}.) \textbf{\hi{लगा}} and \textbf{\hi{लगी}} belong to the verb in this lesson.
\end{quote}

\subsection*{You'll want to know: \hi{लगना}}
\textbf{\hi{लगना}} — \emph{lagnā} — \textquotedblleft{}to seem; to be felt\textquotedblright{}.

\begin{quote}
\textbf{\hi{मुझे} \hi{भूख} \hi{लगी} \hi{है।}}
\end{quote}

\begin{quote}
\emph{mujhe bhūkh lagī hai.} — \textquotedblleft{}I'm hungry.\textquotedblright{}
\end{quote}

Word for word, \textquotedblleft{}hunger has struck me\textquotedblright{}. The one who feels it takes \textbf{\hi{मुझे}}, as in \textbf{\hi{मुझे} \hi{चाय} \hi{चाहिए}}, and the verb follows the thing felt: \textbf{\hi{भूख}} is feminine, so \textbf{\hi{लगी}}.

\begin{quote}
\textbf{\hi{यह} \hi{खाना} \hi{मुझे} \hi{अच्छा} \hi{लगता} \hi{है।}}
\end{quote}

\begin{quote}
\emph{yah khānā mujhe acchā lagtā hai.} — \textquotedblleft{}I like this food.\textquotedblright{}
\end{quote}

Word for word, \textquotedblleft{}this food seems good to me\textquotedblright{}. \textbf{\hi{खाना}} is masculine, so \textbf{\hi{लगता}}. It says much what \textbf{\hi{मुझे} … \hi{पसंद} \hi{है}} says.

\begin{cousinweb}
\textbf{\hi{लगना}} (\emph{lagnā}) is a tadbhava. CDIAL traces it to Sanskrit \emph{lagyati}, \textquotedblleft{}adheres to\textquotedblright{} (CDIAL 10895), and gives \textquotedblleft{}to be attached to\textquotedblright{} as a Hindi sense: a hunger, a hurt or a liking attaches itself to the one who feels it.
\end{cousinweb}

\subsection*{Your turn}
Say these aloud:

\begin{itemize}
\raggedright
  \item \emph{lagnā}, then \emph{mujhe bhūkh lagī hai}
  \item \textquotedblleft{}I'm thirsty\textquotedblright{} — \emph{mujhe pyās lagī hai}
  \item \textquotedblleft{}I like this food\textquotedblright{} — \emph{yah khānā mujhe acchā lagtā hai}
\end{itemize}

\begin{tcolorbox}[breakable,colback=teal!4,colframe=teal!35!black,title={Before you move on}]
What does \textbf{\hi{लगना}} mean? (\textquotedblleft{}To seem, to be felt\textquotedblright{}.) \textquotedblleft{}I'm hungry\textquotedblright{}? (\textbf{\hi{मुझे} \hi{भूख} \hi{लगी} \hi{है}}.) Why \textbf{\hi{लगी}}? (\textbf{\hi{भूख}} is feminine, and the verb follows the thing felt.)
\end{tcolorbox}
\section[\texorpdfstring{\hi{ऐसा} (aisā)}{ऐसा (aisā)}]{\hi{ऐसा} (aisā) — like this, such}

\label{lesson:HI-C280-aisa}



\begin{quote}
Before the new one: ask \textquotedblleft{}How are you?\textquotedblright{} (\textbf{\hi{आप} \hi{कैसे} \hi{हैं}?}) \textbf{\hi{कैसा}} asks \textquotedblleft{}what kind?\textquotedblright{}. This lesson's word answers it.
\end{quote}

\subsection*{You'll want to know: \hi{ऐसा}}
\textbf{\hi{ऐसा}} — \emph{aisā} — \textquotedblleft{}like this, such; so\textquotedblright{}.

\begin{quote}
\textbf{\hi{मुझे} \hi{ऐसा} \hi{लगता} \hi{है।}}
\end{quote}

\begin{quote}
\emph{mujhe aisā lagtā hai.} — \textquotedblleft{}I think so.\textquotedblright{}
\end{quote}

Word for word, \textquotedblleft{}it seems so to me\textquotedblright{}. \textbf{\hi{ऐसा}} agrees, as \textbf{\hi{कैसा}} does: \textbf{\hi{ऐसा}}, \textbf{\hi{ऐसी}}, \textbf{\hi{ऐसे}} (\emph{aisā}, \emph{aisī}, \emph{aise}).

\begin{quote}
\textbf{\hi{ऐसा} \hi{मत} \hi{करो।}}
\end{quote}

\begin{quote}
\emph{aisā mat karo.} — \textquotedblleft{}Don't do that.\textquotedblright{}
\end{quote}

\begin{itemize}
\raggedright
  \item \textbf{\hi{ऐसी} \hi{बात}} (\emph{aisī bāt}) — such a thing; \textbf{\hi{बात}} is feminine
  \item \textbf{\hi{ऐसे} \hi{लिखिए।}} (\emph{aise likhiye.}) — Write it like this.
\end{itemize}

Its first letter is new to you: \textbf{\hi{ऐ}}. The next lesson is about it.

\begin{cousinweb}
\textbf{\hi{ऐसा}} (\emph{aisā}) is a tadbhava, worn down from Sanskrit \emph{īdṛśa}, \textquotedblleft{}such, like this\textquotedblright{} (CDIAL 1611). Its question partner has the same history: \textbf{\hi{कैसा}} goes back to \emph{kīdṛśa}, \textquotedblleft{}of what kind?\textquotedblright{} (CDIAL 3197).
\end{cousinweb}

\subsection*{Your turn}
Say these aloud:

\begin{itemize}
\raggedright
  \item \emph{aisā}, \emph{aisī}, \emph{aise}
  \item \textquotedblleft{}I think so\textquotedblright{} — \emph{mujhe aisā lagtā hai}
  \item \textquotedblleft{}Don't do that\textquotedblright{} — \emph{aisā mat karo}
\end{itemize}

\begin{tcolorbox}[breakable,colback=teal!4,colframe=teal!35!black,title={Before you move on}]
What does \textbf{\hi{ऐसा}} mean? (\textquotedblleft{}Like this, such\textquotedblright{}.) Its question partner? (\textbf{\hi{कैसा}}.) \textquotedblleft{}I think so\textquotedblright{}? (\textbf{\hi{मुझे} \hi{ऐसा} \hi{लगता} \hi{है}}.)
\end{tcolorbox}
\section[\texorpdfstring{\hi{ऐ} (ai)}{ऐ (ai)}]{\hi{ऐ} — the vowel you have read since \hi{है}, standing alone}

\label{lesson:HI-W280-letter-ai}



\begin{quote}
Before the new one: say \textbf{\hi{ऐसा}}, \textquotedblleft{}like this\textquotedblright{}. Then look at \textbf{\hi{है}}. How many strokes sit above its head-line? (Two: \textbf{\hi{◌ै}}, \emph{ai}.)

That sign has a \textbf{standing form}, and \textbf{\hi{ऐसा}} opens with it.
\end{quote}

\subsection*{Script you'll notice: \hi{ऐ}}
\textbf{\hi{ऐ}} — \emph{ai}.

It is the \textbf{independent vowel}. \textbf{\hi{◌ै}} is how \emph{ai} looks when it rides on a consonant (\textbf{\hi{है}}, \textbf{\hi{मैं}}, \textbf{\hi{पैसा}}); \textbf{\hi{ऐ}} is how it looks when it stands on its own, with no consonant to carry it.

What it is made of:

\begin{itemize}
\raggedright
  \item the long stem and descending tail of \textbf{\hi{ए}}
  \item the shorter hooked stem of \textbf{\hi{ए}}
  \item a separate upper arc
  \item the top shirorekhā
\end{itemize}

In print, then, \textbf{\hi{ऐ}} is \textbf{\hi{ए}} with one arc that rises from the right end of the head-line and curves up and over to the left.

\subsection*{Script: two shapes, one sound}
\noindent
\begin{tabularx}{\linewidth}{@{}>{\raggedright\arraybackslash}X>{\raggedright\arraybackslash}X@{}}
\toprule
\textbf{what you see} & \textbf{where it stands} \\
\midrule
\hi{◌ै} & on a consonant: \textbf{\hi{है}}, \textbf{\hi{मैं}} \\
\hi{ऐ} & alone: \textbf{\hi{ऐसा}} \\
\bottomrule
\end{tabularx}

Count the strokes above the line. \textbf{\hi{◌े}} has one and \textbf{\hi{◌ै}} two; \textbf{\hi{ए}} has none and \textbf{\hi{ऐ}} one. Miscount, and \emph{e} becomes \emph{ai}, or the other way round.

\subsection*{Writing: \hi{ऐ}}
\hlblockfigure{\detokenize{figures/HI-W280-letter-ai-filmstrip.pdf}}{How \hi{ऐ} is written, stroke by stroke}

\begin{itemize}
\raggedright
  \item \textbf{1.} descend the long left stem, curve right through the lower shoulder, and sweep down through the tail without lifting
  \item \textbf{2.} lift and descend the shorter right stem into its inward hook
  \item \textbf{3.} lift at the headline junction and sweep the upper arc upward and left
  \item \textbf{4.} lift and draw the top shirorekhā left-to-right
\end{itemize}

\textbf{Pen lifts: 3.} The first two strokes are the ones you write for \textbf{\hi{ए}}; the arc comes third, and the head-line last.

\begin{quote}
Verified modern printed teaching form; everyday handwriting may simplify the curves.
\end{quote}

\begin{quote}
This is one attested teaching order and not a national standard — handwriting here is taught with school-to-school variation. Source: Saurmandal, ‘Devanagari \hi{ऐ} stroke order.svg’, panels 1–4, Wikimedia Commons, 5 August 2023.
\end{quote}

\subsection*{Your turn}
\begin{itemize}
\raggedright
  \item \emph{Look:} at the sign and the letter side by side, and say what they share
\end{itemize}

\begin{quote}
\hi{◌ै}  ·  \hi{ऐ}
\end{quote}

\begin{itemize}
\raggedright
  \item \emph{Trace it:} \hi{ऐ} three times, saying \emph{ai} as you finish each one
  \item \emph{Write it:} \textbf{\hi{ऐसा}}, then \textbf{\hi{ऐसी}}
\end{itemize}

\begin{tcolorbox}[breakable,colback=teal!4,colframe=teal!35!black,title={Before you move on}]
What sound is \textbf{\hi{ऐ}}? (\emph{ai}.) What is its sitting form? (\textbf{\hi{◌ै}}, as in \textbf{\hi{है}}.) How is \textbf{\hi{ऐ}} unlike \textbf{\hi{ए}}? (One arc rises above its head-line.)
\end{tcolorbox}
\section[\texorpdfstring{\hi{मुझे} \hi{लगता} \hi{है} \hi{कि} … (mujhe lagtā hai ki …)}{मुझे लगता है कि … (mujhe lagtā hai ki …)}]{\hi{मुझे} \hi{लगता} \hi{है} \hi{कि} … (mujhe lagtā hai ki …) — what I think}

\label{lesson:HI-C280-mujhelagtahai}



\begin{quote}
Say \textquotedblleft{}I think so\textquotedblright{}. (\textbf{\hi{मुझे} \hi{ऐसा} \hi{लगता} \hi{है}}.) Then \textbf{\hi{कि}}, \textquotedblleft{}that\textquotedblright{}, the word that turns a sentence into the thing you think or know.
\end{quote}

\subsection*{You'll want to know: \hi{मुझे} \hi{लगता} \hi{है} \hi{कि}}
\begin{quote}
\textbf{\hi{मुझे} \hi{लगता} \hi{है} \hi{कि} \hi{आज} \hi{बारिश} \hi{होगी।}}
\end{quote}

\begin{quote}
\emph{mujhe lagtā hai ki āj bāriś hogī.} — \textquotedblleft{}I think it will rain today.\textquotedblright{}
\end{quote}

Word for word, \textquotedblleft{}it seems to me that …\textquotedblright{}. The thinker takes \textbf{\hi{मुझे}}, and \textbf{\hi{लगता} \hi{है}} stays as it is whatever follows \textbf{\hi{कि}}.

\begin{itemize}
\raggedright
  \item \textbf{\hi{आपको} \hi{क्या} \hi{लगता} \hi{है}?} (\emph{āpko kyā lagtā hai?}) — What do you think?
  \item \textbf{\hi{तुम्हें} \hi{क्या} \hi{लगता} \hi{है}?} (\emph{tumhẽ kyā lagtā hai?}) — the same, to a friend
  \item \textbf{\hi{मुझे} \hi{नहीं} \hi{लगता} \hi{कि} \hi{वह} \hi{आएगा।}} (\emph{mujhe nahīṁ lagtā ki vah āegā.}) — I don't think he will come. With \textbf{\hi{नहीं}}, the \textbf{\hi{है}} drops.
  \item \textbf{\hi{मुझे} \hi{लगा} \hi{कि} \hi{तुम} \hi{नहीं} \hi{आओगे।}} (\emph{mujhe lagā ki tum nahīṁ āoge.}) — I thought you wouldn't come.
\end{itemize}

\subsection*{You'll want to know: in my opinion}
\begin{quote}
\textbf{\hi{मेरी} \hi{राय} \hi{में} \hi{हिंदी} \hi{आसान} \hi{है।}}
\end{quote}

\begin{quote}
\emph{merī rāy meṁ hindī āsān hai.} — \textquotedblleft{}In my opinion Hindi is easy.\textquotedblright{}
\end{quote}

\textbf{\hi{राय}} is feminine, so \textbf{\hi{मेरी} \hi{राय} \hi{में}}, \textquotedblleft{}in my opinion\textquotedblright{}. \textbf{\hi{विचार}} is masculine: \textbf{\hi{मेरे} \hi{विचार} \hi{से}} (\emph{mere vicār se}) says the same, a little more formally. To ask: \textbf{\hi{आपकी} \hi{क्या} \hi{राय} \hi{है}?} (\emph{āpkī kyā rāy hai?}), \textquotedblleft{}What is your opinion?\textquotedblright{}

\textbf{\hi{मैं} \hi{सोचता} \hi{हूँ} \hi{कि} …} (\emph{maiṁ soctā hūṁ ki …}) puts the thinker in charge: it is for a thought you are turning over, not for an impression.

\subsection*{Your turn}
Say these aloud:

\begin{itemize}
\raggedright
  \item \textquotedblleft{}What do you think?\textquotedblright{} — \emph{āpko kyā lagtā hai?}
  \item \textquotedblleft{}I don't think he will come\textquotedblright{} — \emph{mujhe nahīṁ lagtā ki vah āegā}
  \item \textquotedblleft{}I thought you wouldn't come\textquotedblright{} — \emph{mujhe lagā ki tum nahīṁ āoge}
\end{itemize}

\begin{tcolorbox}[breakable,colback=teal!4,colframe=teal!35!black,title={Before you move on}]
\textquotedblleft{}I think that …\textquotedblright{}? (\textbf{\hi{मुझे} \hi{लगता} \hi{है} \hi{कि} …}) \textquotedblleft{}What do you think?\textquotedblright{} (\textbf{\hi{आपको} \hi{क्या} \hi{लगता} \hi{है}?}) \textquotedblleft{}In my opinion\textquotedblright{}? (\textbf{\hi{मेरी} \hi{राय} \hi{में}}: \textbf{\hi{राय}} is feminine.) \textquotedblleft{}I thought that …\textquotedblright{}? (\textbf{\hi{मुझे} \hi{लगा} \hi{कि} …})
\end{tcolorbox}
\section[\texorpdfstring{\hi{की} \hi{वजह} \hi{से} (kī vajah se)}{की वजह से (kī vajah se)}]{\hi{की} \hi{वजह} \hi{से} (kī vajah se) — because of}

\label{lesson:HI-C280-kivajahse}



\begin{quote}
Answer a \textbf{\hi{क्यों}} with \textbf{\hi{क्योंकि}}: \textquotedblleft{}Why didn't you come? Because I was ill.\textquotedblright{} (\textbf{\hi{तुम} \hi{क्यों} \hi{नहीं} \hi{आए}? \hi{क्योंकि} \hi{मैं} \hi{बीमार} \hi{था।}})

\textbf{\hi{क्योंकि}} needs a whole sentence after it. This lesson gives a reason in a single noun.
\end{quote}

\subsection*{You'll want to know: \hi{की} \hi{वजह} \hi{से}, \hi{के} \hi{कारण}}
\begin{quote}
\textbf{\hi{बारिश} \hi{की} \hi{वजह} \hi{से} \hi{हम} \hi{घर} \hi{पर} \hi{रहे।}}
\end{quote}

\begin{quote}
\emph{bāriś kī vajah se ham ghar par rahe.} — \textquotedblleft{}Because of the rain we stayed at home.\textquotedblright{}
\end{quote}

\textbf{\hi{वजह}} (\emph{vajah}), \textquotedblleft{}a reason\textquotedblright{}, is feminine, so \textbf{\hi{की} \hi{वजह} \hi{से}}: \textquotedblleft{}from the reason of\textquotedblright{}. \textbf{\hi{कारण}}, which you know, is masculine and does the same job a little more formally: \textbf{\hi{के} \hi{कारण}} (\emph{ke kāraṇ}).

\begin{itemize}
\raggedright
  \item \textbf{\hi{भीड़} \hi{की} \hi{वजह} \hi{से} \hi{मुझे} \hi{देर} \hi{हुई।}} (\emph{bhīṛ kī vajah se mujhe der huī.}) — Because of the crowd I was late.
  \item \textbf{\hi{बीमारी} \hi{के} \hi{कारण} \hi{वह} \hi{नहीं} \hi{आया।}} (\emph{bīmārī ke kāraṇ vah nahīṁ āyā.}) — Because of illness he didn't come.
\end{itemize}

A pronoun takes its \textbf{-\hi{की}} form: \textbf{\hi{मेरी} \hi{वजह} \hi{से}} (\emph{merī vajah se}), \textquotedblleft{}because of me\textquotedblright{}; \textbf{\hi{उसकी} \hi{वजह} \hi{से}} (\emph{uskī vajah se}), \textquotedblleft{}because of him\textquotedblright{}, or her.

\subsection*{You'll want to know: the reason is that …}
\begin{quote}
\textbf{\hi{वजह} \hi{यह} \hi{है} \hi{कि} \hi{मैं} \hi{बीमार} \hi{था।}}
\end{quote}

\begin{quote}
\emph{vajah yah hai ki maiṁ bīmār thā.} — \textquotedblleft{}The reason is that I was ill.\textquotedblright{}
\end{quote}

\textbf{\hi{वजह} \hi{यह} \hi{है} \hi{कि} …} names the reason and then gives it as a whole sentence. Mind the order of the three ways:

\noindent
\begin{tabularx}{\linewidth}{@{}>{\raggedright\arraybackslash}X>{\raggedright\arraybackslash}X@{}}
\toprule
\textbf{the reason} & \textbf{where it stands} \\
\midrule
\textbf{\hi{बारिश} \hi{की} \hi{वजह} \hi{से}} & before the result \\
\textbf{\hi{वजह} \hi{यह} \hi{है} \hi{कि} …} & in a sentence of its own, after the result \\
\textbf{\hi{क्योंकि} …} & after the result, in the same sentence \\
\bottomrule
\end{tabularx}

\subsection*{Your turn}
Say these aloud:

\begin{itemize}
\raggedright
  \item \textquotedblleft{}Because of the crowd I was late\textquotedblright{} — \emph{bhīṛ kī vajah se mujhe der huī}
  \item \textquotedblleft{}He didn't come, because he was ill\textquotedblright{} — \emph{vah nahīṁ āyā, kyoṁki vah bīmār thā}
  \item \textquotedblleft{}Because of me she was late\textquotedblright{} — \emph{merī vajah se use der huī}
\end{itemize}

\begin{tcolorbox}[breakable,colback=teal!4,colframe=teal!35!black,title={Before you move on}]
\textquotedblleft{}Because of the rain\textquotedblright{}? (\textbf{\hi{बारिश} \hi{की} \hi{वजह} \hi{से}}.) Why \textbf{\hi{की}}? (\textbf{\hi{वजह}} is feminine.) Its formal partner? (\textbf{\hi{के} \hi{कारण}}.) \textquotedblleft{}The reason is that …\textquotedblright{}? (\textbf{\hi{वजह} \hi{यह} \hi{है} \hi{कि} …}) Which one comes after the result, in the same sentence? (\textbf{\hi{क्योंकि}}.)
\end{tcolorbox}
\section[(review)]{Review — what I think, and why}

\label{lesson:HI-R280-opinion-and-reason}



\begin{quote}
Say \textquotedblleft{}I'm hungry\textquotedblright{} and \textquotedblleft{}I think so\textquotedblright{}. (\emph{mujhe bhūkh lagī hai}, \emph{mujhe aisā lagtā hai}.) Then write \textbf{\hi{ऐ}} from memory: the strokes of \textbf{\hi{ए}}, the arc, the head-line.
\end{quote}

\subsection*{Your turn}
\begin{itemize}
\raggedright
  \item \emph{Say it:} I think it will rain tomorrow.
  \item \emph{Say it:} What do you think? In my opinion this book is good.
  \item \emph{Say it:} Because of the rain the train came late.
  \item \emph{Say it:} I'm thirsty. I like this tea.
  \item \emph{Write it:} \textbf{\hi{ऐसा}}, then \textbf{\hi{मुझे} \hi{ऐसा} \hi{लगता} \hi{है।}}
\end{itemize}

\begin{tcolorbox}[breakable,colback=teal!4,colframe=teal!35!black,title={Before you move on}]
\textquotedblleft{}I think that …\textquotedblright{}? (\textbf{\hi{मुझे} \hi{लगता} \hi{है} \hi{कि} …}) \textquotedblleft{}In my opinion\textquotedblright{}? (\textbf{\hi{मेरी} \hi{राय} \hi{में}}.) \textquotedblleft{}Because of\textquotedblright{}? (\textbf{\hi{की} \hi{वजह} \hi{से}}, or \textbf{\hi{के} \hi{कारण}}.) \textquotedblleft{}The reason is that …\textquotedblright{}? (\textbf{\hi{वजह} \hi{यह} \hi{है} \hi{कि} …}) How is \textbf{\hi{ऐ}} unlike \textbf{\hi{ए}}? (One arc above the head-line.)
\end{tcolorbox}
